% =============================================================================
% MATERIALS AND METHODS
% Status: DRAFT - Core methodology documented
% =============================================================================

\subsection{Study Area and Data Collection}

\subsubsection{Location}
Santa Ines agricultural field, Chile. Tomato cultivation site monitored via drone surveys.

\todo{Add exact coordinates, field size, crop details}

\subsubsection{Drone Imagery Dataset}
\begin{itemize}
    \item \textbf{Platform}: DJI drone (model TBD)
    \item \textbf{Image count}: 97 images
    \item \textbf{File format}: JPEG, approximately 10-12 MB per image
    \item \textbf{Resolution}: High-resolution (dimensions TBD from metadata)
    \item \textbf{Date}: January 2026
\end{itemize}

\subsection{Software Environment}

\subsubsection{Dataset Management}
FiftyOne v1.12.0 \citep{fiftyone} was used for dataset management, embedding computation, and visualization. FiftyOne provides:
\begin{itemize}
    \item Efficient handling of large image datasets
    \item Built-in integration with embedding models
    \item Interactive visualization tools
    \item Similarity search functionality
\end{itemize}

\subsubsection{Analysis Tools}
\begin{itemize}
    \item \textbf{Python 3.12}: Data processing and embedding computation
    \item \textbf{R 4.3.3}: Statistical analysis and visualization (ggplot2)
    \item \textbf{PyTorch}: Deep learning backend for CLIP model
\end{itemize}

\subsection{Embedding Computation}

\subsubsection{Model Selection}
CLIP ViT-B/32 (Vision Transformer, Base, 32x32 patches) was selected for embedding computation. This model:
\begin{itemize}
    \item Provides 512-dimensional embeddings per image
    \item Was pretrained on 400M image-text pairs
    \item Offers good balance of computational efficiency and semantic richness
\end{itemize}

\subsubsection{Dimensionality Reduction}
UMAP (Uniform Manifold Approximation and Projection) was applied to reduce embeddings from 512D to 2D for visualization:

\begin{equation}
    \mathbf{z}_i = \text{UMAP}(\mathbf{e}_i), \quad \mathbf{e}_i \in \mathbb{R}^{512}, \quad \mathbf{z}_i \in \mathbb{R}^{2}
\end{equation}

\todo{Add UMAP hyperparameters: n\_neighbors, min\_dist, metric}

\subsubsection{Implementation}
\begin{lstlisting}[language=Python, caption=Embedding computation with FiftyOne]
import fiftyone as fo
import fiftyone.brain as fob

dataset = fo.load_dataset("santa_ines_dron")

# Compute CLIP embeddings with UMAP projection
fob.compute_visualization(
    dataset,
    brain_key="clip_viz",
    model="clip-vit-base32-torch",
    method="umap"
)

# Compute similarity index for retrieval
fob.compute_similarity(
    dataset,
    brain_key="clip_similarity",
    embeddings="clip_viz"
)
\end{lstlisting}

\subsection{Clustering and Classification}

Image clusters were identified through visual inspection of the UMAP projection. Two distinct clusters emerged:

\begin{enumerate}
    \item \textbf{Main cluster} (89 images): Concentrated in UMAP region $x \in [6.8, 10.5]$, $y \in [0.5, 4.5]$
    \item \textbf{Outlier cluster} (8 images): Isolated at $x \approx 11.5$, $y \approx -0.7$
\end{enumerate}

Images were tagged based on cluster membership for subsequent analysis.

\subsection{Similarity Network Analysis}

To quantify cluster structure, a k-nearest neighbor graph was constructed:

\begin{itemize}
    \item Each image connected to its $k=3$ nearest neighbors in embedding space
    \item Distance metric: Euclidean distance in 2D UMAP space
    \item Edge classification: within-category vs. cross-category connections
\end{itemize}

\subsection{Statistical Analysis}

Cluster cohesion was evaluated using:
\begin{itemize}
    \item Mean distance to $k$ nearest neighbors
    \item Inter-cluster centroid distance
    \item Proportion of cross-category edges in similarity network
\end{itemize}

All statistical analyses and visualizations were performed in R using ggplot2.
